\scanpage{8}
\mspage{7}

La commutation aux $\varprojlim$ adiques peut aussi s'exprimer ainsi,

($L$) Pour tt \deleted{\unsure{schéma}} $T$ \inserted{sur $S$} spectre d'un anneau local noeth. \deleted{complet sur $S$}\inserted{(qu'on peut supposer complet, à corps rés. alg. clos.)}, et tel que $Y_T$ majore un voisinage \unsure{ouvert induit|ouvert infinit.} de la fibre spéciale de $X_T$, on a \deleted{$\unsure{X_T}$} $X_T=Y_T$.

\uline{Corollaire} Cas $X$ propre sur $S$; \uuline{O.K.}

Si $X$ est un groupe, $Y$ un sous-groupe, on sait que $Y_T$ \deleted{contient}\inserted{majore} un voisinage \unsure{ouvert induit|ouvert infinit.} de la section neutre, donc $Y_T-X_{\unsure{x|T}}$ est \inserted{au \uline{moins}} \unsure{étale}, \deleted{donc} étant une imm \deleted{fermée}, est une immersion ouverte \deleted{et fermée}. \deleted{\unsure{Donc dans ce cas}}

\inserted{De plus, la vérification $Y_T=X_T$ \unsure{est de|est la} \deleted{\unsure{nécessaire}} \unsure{très|par} ``\unsure{localité|constructibilité}'', donc se ramène au cas où $T$ est le spectre d'un anneau de valuation discrète. Donc}

\uline{Corollaire} Dans ce cas, la condition envisagée équivaut à la suivante: Pour tt $T$ sur $S$ spectre d'un anneau \deleted{de valuation}\deleted{local noeth.}\inserted{\uline{de valuation discrète}} (qu'on peut supposer complet à corps rés. alg. clos), tel que $Y_T$ soit un sous-groupe ouvert et fermé de $X_T$ qui contient la fibre \deleted{générique}\inserted{spéciale}, on a $Y_T=X_T$. Cette condition est vérifiée dans chacun des cas suivants.

a) $X$ à \uline{fibres connexes}

b) Il \uline{existe} une \uline{famille de $S$-\deleted{groupes finis}\inserted{schémas}} $\unsure{G_i^?|G_i^!}$ de $G$, finis sur $S$, et tels que pour chaque $s\in S$, la réunion ensembliste des $G_{i\bar s}.G_{\bar s}^0$ soit $G_{\bar s}$.
